\clearpage\begingroup\NoSourceHead\SourceLeaf{460}
\renewcommand{\thefootnote}{\arabic{footnote})}\setcounter{footnote}{0}
\noindent\hspace*{\parindent}Die Beziehungen von \textsc{Gauss} zur Variationsrechnung finden sich im wesentlichen
in den beiden Arbeiten: »\textit{Principia generalia theoriae figurae fluidorum in
statu aequilibrii}«\footnote{Werke, Bd. V, S. 29.} und »\textit{Disquisitiones generales circa superficies curvas}«\footnote{Werke, Bd. IV, S. 217.}. In
der ersten Arbeit hat \textsc{Gauss} \textls[100]{unmittelbar} an einer wichtigen Stelle in die
Entwicklung der Variationsrechnung eingegriffen; in der zweiten spielt die
Variationsrechnung zwar unmittelbar nur eine untergeordnete Rolle, aber um
so grösser ist ihre \textls[100]{mittelbare} Wirkung auf die Variationsrechnung gewesen,
insofern die \textsc{Gauss}schen Untersuchungen über geodätische Linien aufs engste
mit den fundamentalsten Fragestellungen der neueren Variationsrechnung in
Zusammenhang stehen und dazu vielfach den Anstoss gegeben haben. Ferner
findet sich eine wichtige Anwendung der Variationsrechnung in der Arbeit:
»\textit{Allgemeine Lehrsätze in Beziehung auf die im verkehrten Verhältnisse des Quadrats
der Entfernung wirkenden Anziehungs- und Abstossungskräfte}«\footnote{Werke, Bd. V, S. 195.}, in der \textsc{Gauss}
ein Existenztheorem der Potentialtheorie nach einer Methode beweist, deren
Grundgedanke mit dem \textsc{Dirichlet}schen Prinzip identisch ist. Mit diesen drei
Arbeiten beschäftigen sich der Reihe nach die drei ersten Teile der folgenden
Darstellung. Daran schliesst sich dann noch ein vierter Teil, in dem verschiedene
vereinzelte kürzere Anwendungen besprochen werden, die \textsc{Gauss}
gelegentlich von der Variationsrechnung gemacht hat, und die sich in seinem
Nachlass gefunden haben.
\par\smallskip\noindent\hfill{\small 1\textsuperscript{*}}\par
\clearpage\endgroup
